\section{Introduction and conventions}
\label{sec:intro}

\subsection{What the code solves}
\label{sec:intro:what}

EXHALE integrates the one-dimensional, spherically symmetric equations of
radiation hydrodynamics for the escaping upper atmosphere of an irradiated
planet.  The radial domain runs from a base near the lower atmosphere, which
a run may place at the planet radius or at a stated pressure level (the
$1\,\mu$bar level of the \citet{Koskinen2022} handoff is the one the molecular
configurations use), out to the Roche-lobe (L1) radius of the planet--star
system or, in the spherical domain mode, to a radius the user states.  The
unknowns advanced in time are the conserved hydrodynamic state
$u = (\rho,\,\rho v,\,\mathcal{E})$ on a finite-volume radial grid, together with the
ionization and chemical state of the gas: hydrogen and helium always, the
metastable $\mathrm{He}\,2^3S$ level, a set of trace metals, and a molecular
network carrying $\mathrm{H}_2$, $\mathrm{H}_2^+$, $\mathrm{H}_3^+$,
$\mathrm{HeH}^+$ and, optionally, the oxygen carriers OH, $\mathrm{H_2O}$ and
CO.  The radiative source term couples the two: the stellar spectrum is
attenuated along the radial column, photoionizes and heats the gas, and the
gas cools through recombination, collisional and line processes whose rates
depend on the composition the ionization solve returns.

The code is a heavily extended fork of ATES v2.0 \citep{Caldiroli2021,
Biassoni2024}.  The hydrodynamic core, the finite-volume discretization, the
three-stage strong-stability-preserving Runge-Kutta integration and the
choice of approximate Riemann solvers are recognizably those of ATES; the
equation of state, the lower boundary condition, the reconstruction, the
steady-state solvers, the transported species and essentially all of the
chemistry are not.  Every physics extension is selected at run time by the
presence of a key in \texttt{input.inp} or of a companion file
(\texttt{metals.inp}, \texttt{opacity.inp}, \texttt{base.inp}), and every one
of them defaults to off, so that a configuration that states none of them
reproduces the atomic H/He problem.

\subsection{Units}
\label{sec:intro:units}

The equations are solved in dimensionless form.  The normalization is fixed
once, in \texttt{input\_read.f90}, from the planet radius $R_0$, the base
temperature $T_0$, and the base number density $n_0$:
%
\begin{equation}
  v_0 = \sqrt{\frac{k_B T_0}{m_{\rm H}}}, \qquad
  t_s = \frac{R_0}{v_0}, \qquad
  p_0 = n_0\,m_{\rm H}\,v_0^2, \qquad
  q_0 = \frac{n_0\,m_{\rm H}\,v_0^3}{R_0},
  \label{eq:units}
\end{equation}
%
where $m_{\rm H} = 1.67353284\times10^{-24}$~g is the mass of the hydrogen
\emph{atom}, not the atomic mass unit.  Radii are in units of $R_0$, so the
base level sits at $r = 1$; velocities in $v_0$; mass densities in
$n_0m_{\rm H}$; pressures in $p_0$; temperatures in $T_0$; and volumetric energy
rates in $q_0$.  The mass unit being the hydrogen atom, the density is the
number of hydrogen-atom masses per unit volume, and a species contributes
$m_{\rm species}/m_{\rm H}$ to it; the species table therefore carries
$m(\mathrm{He\,I}) = 3.9715259$ and not $4$.

The gravitational potential is measured in $v_0^2$, so it is fixed by the
single dimensionless group
%
\begin{equation}
  b_0 = \frac{G M_p m_{\rm H}}{k_B T_0 R_0},
  \label{eq:b0}
\end{equation}
%
the Jeans escape parameter at the base.  Only three further dimensionless
numbers enter the potential and the domain: $\tilde a = a_{\rm orb}/R_0$, the
mass ratio $M_\star/M_p$, and the outer radius.

The base level itself may be stated in three ways, and a run that states two
that disagree is stopped at startup: \texttt{Log10 lower boundary number
density} gives $n_0$ directly; \texttt{Base BC: pressure} states the base
pressure in microbar and $n_0$ follows from
$n_0 = p_{\rm base}/(k_B T_0\,\hat n_{\rm bc})$; and a lower-atmosphere
handoff (\texttt{base.inp} or a profile file) states $p_{\rm base}$ in bar
and $n_0$ follows the same way.  Here $\hat n_{\rm bc}$ is the base particle
count per H+He nucleus, which the composition module derives from the
element ratios and, with a molecular base, from the H nuclei bound into
$\mathrm{H}_2$ (Sec.~\ref{sec:hydro:eos}).

\subsection{Notation}
\label{sec:intro:notation}

One symbol stands for one quantity throughout this document.  The table below
fixes the symbols that appear in more than one section.  Symbols introduced
inside a published fit and defined where they appear (the $x$ of a
Verner-type cross section, the $\lambda$ of a recombination-cooling fit, the
$\alpha$ of a spectral index) are local to that fit and are not in the table.

\begin{center}
\small
\begin{tabular}{@{}>{\raggedright\arraybackslash}p{0.13\textwidth}>{\raggedright\arraybackslash}p{0.46\textwidth}>{\raggedright\arraybackslash}p{0.33\textwidth}@{}}
\toprule
symbol & quantity & note \\
\midrule
$u$ & conserved state $(\rho,\rho v,\mathcal{E})$ &
  the unknown of the hydrodynamics, Sec.~\ref{sec:hydro:euler} \\
$W$ & primitive state $(\rho,v,p)$ & what the reconstruction works on \\
$\rho$, $v$, $p$, $T$ & mass density, radial velocity, pressure,
  temperature & \\
$\mathcal{E}$, $\rho e$, $e$, $\varepsilon$ & total energy density, internal
  energy density, internal energy per unit mass, internal energy per
  particle & $\mathcal{E} = \tfrac12\rho v^2+\rho e$ \\
$\Phi$ & gravitational potential & $\Phi_i$, $\Phi_c$ at interfaces and cell
  centers \\
$\gamma_{\rm eff}$, $c_s$ & effective adiabatic index, sound speed &
  from the caloric equation of state \eqref{eq:caloric} \\
$m_{\rm H}$ & mass of the hydrogen atom & the code's mass unit,
  \eqref{eq:units} \\
$\mu$, $\kappa$ & dynamic viscosity, heat conductivity &
  \eqref{eq:transport_coeff} only \\
$n_X$, $n_{\rm tot}$, $n_e$ & number density of species $X$, of all heavy
  particles, of free electrons & \\
$x_{X^{q+}}$, $f_i$ & ionization fraction of a stage, mixing ratio
  $n_i/n_{\rm tot}$ & fractions of the element's own nuclei \\
$E$, $h\nu$ & photon energy & always subscripted where a particular
  energy is meant ($E_k$, $E_{\rm th}$) \\
$F_E$, $\mathcal{F}_k$ & incident and attenuated flux per unit photon
  energy & \eqref{eq:rad:workingflux} \\
$\sigma$, $\tau$ & photoionization cross section, optical depth & \\
$\Gamma$, $\alpha$, $\beta$ & photoionization rate, recombination
  coefficient, electron-impact ionization coefficient &
  Sec.~\ref{sec:ionization:cellsystem} \\
$\mathcal{H}$, $\Lambda$ & volumetric heating rate, volumetric cooling
  rate & erg\,cm$^{-3}$\,s$^{-1}$; $Q = \mathcal{H}-\Lambda$ \\
$\xi_{\rm day}$ & dayside dilution of every stellar beam &
  Table~\ref{tab:2d} \\
$J_i$, $J$ & diffusive number flux of a carrier, diffusive mass flux of an
  element & Secs.~\ref{sec:mol:carrier}, \ref{sec:mol:diffusion} \\
$dF$, $S$, $R$ & discrete flux difference, source, stationary residual &
  \eqref{eq:dF1}--\eqref{eq:source}, \eqref{eq:residual} \\
$A_\pm$, $\Delta V_j$, $\Delta r_j$ & face areas, cell volume, cell width &
  \eqref{eq:geom} \\
\bottomrule
\end{tabular}
\end{center}

\subsection{How to read this document}
\label{sec:intro:layout}

Every equation below is the discrete operator the code applies, taken from
the source and not from a derivation.  Where the code's expression differs
from the textbook form (the spherical pressure term under PLM, the
gravitational work term folded into the energy flux, the mean molecular
weight that is never formed), the code's form is given and the difference is
stated.  Routines are named with the file they live in, for instance
\texttt{RK\_rhs.f90}, so that a statement here can be checked against one
place in the tree.

Section~\ref{sec:hydro} states the conservation laws, the equation of state
and the source terms.  The sections that follow it state the radiation field,
the ionization and chemical network, the thermal source and the molecular and
metal chemistry.  Section~\ref{sec:numerics} then gives the grid, the
reconstruction, the Riemann solvers, the time integration and the operator
splitting; Sec.~\ref{sec:bcic} the boundary and initial conditions; and
Sec.~\ref{sec:steady} the stationary residual, its acceptance test and the
solvers that drive it to zero.  A section on the post-processing tools closes
the physics account, and Sec.~\ref{sec:keys} tabulates the input keys that
select physics, with the default each of them takes when the key is absent.
